\chapter{Literature Review}
\label{ch:literature_review}

\section{Introduction}
\label{sec:lit_intro}

The purpose of this chapter is to conduct an extensive, critical, and structured review of the theoretical, regulatory, and empirical literature pertinent to the valuation of equity securities and the value relevance of financial accounting information, with specific focus on commercial banking institutions undergoing structural consolidation.

This chapter establishes a direct link with the research problem and objectives articulated in Chapter~I. It is organized into four major substantive sections:
\begin{itemize}
\item \textbf{Section~\ref{sec:theoretical_review}: Theoretical Review} examines the foundational economic and financial accounting theories that govern equity valuation, residual income dynamics, post-merger information environments, market efficiency, signaling, agency conflicts, and the econometric hazards of price-derived valuation multiples.
\item \textbf{Section~\ref{sec:regulatory_framework}: Review of Legal, Regulatory, and Policy Framework in Nepal} reviews the institutional architecture governing the Nepalese banking sector and capital market, including the regulatory mandates of Nepal Rastra Bank (NRB), the statutory reporting obligations under the Banks and Financial Institutions Act (BAFIA), 2017, and the policy evolution of bank merger directives and capital expansion mandates.
\item \textbf{Section~\ref{sec:empirical_review}: Empirical Review} provides a comprehensive, comparative analysis of prior empirical studies, divided into international investigations across developed and emerging capital markets, foundational Nepalese banking studies, recent investigations of share price determinants (2021--2024), and contemporary literature spanning the early-2023 consolidation wave (2024--2026), synthesized through a detailed summary matrix.
\item \textbf{Section~\ref{sec:research_gap}: Research Gap} synthesizes the critical methodological, econometric, and conceptual shortcomings of the existing literature and explicitly articulates how the present dissertation addresses each identified gap.
\end{itemize}

\section{Theoretical Review}
\label{sec:theoretical_review}

\subsection{The Residual Income Valuation (RIV) Framework}
\label{subsec:ohlson_framework}

In financial economics and accounting theory, equity valuation has evolved from dividend-discounting formulations into models grounded in fundamental accounting numbers. The seminal theoretical breakthrough linking financial accounting disclosures directly to equity market values was established by \citet{ohlson1995} and extended by \citet{feltham1995}. The Residual Income Valuation (RIV) framework---often designated as the Edwards-Bell-Ohlson (EBO) model---reconceptualizes firm value by expressing the market value of common equity as the sum of currently recognized book value of equity and the present value of expected future residual income (abnormal earnings).

The mathematical foundation of the Ohlson framework rests on four interrelated structural building blocks:

\begin{enumerate}[label=\textbf{\arabic*.}]
\item \textbf{The Present Value of Expected Dividends (PVED):}
Consistent with neoclassical asset pricing theory, the intrinsic market value of a firm's common equity at date $t$ ($P_t$) equals the discounted present value of all expected future net cash distributions (dividends, $D_{t+\tau}$) paid to equity holders:
\begin{equation}
P_t = \sum_{\tau=1}^{\infty} \frac{\mathbb{E}_t[D_{t+\tau}]}{(1 + r)^\tau}
\end{equation}
where $\mathbb{E}_t[\cdot]$ denotes the expectations operator conditioned on information available at date $t$, and $r$ represents the constant cost of equity capital (assumed strictly positive).

\item \textbf{The Clean Surplus Relation (CSR):}
The Ohlson model imposes the fundamental accounting condition of clean surplus accounting. Under CSR, all changes in the book value of equity ($BV_t$) between dates $t-1$ and $t$ must flow through the income statement as net income ($NI_t$), with the sole exception of transactions with equity holders in their capacity as owners (dividends paid, $D_t$, and capital contributions/withdrawals):
\begin{equation}
BV_t = BV_{t-1} + NI_t - D_t
\end{equation}
or equivalently:
\begin{equation}
D_t = NI_t - (BV_t - BV_{t-1})
\end{equation}
The clean surplus relation ensures accounting consistency by preventing firms from bypassing the statement of profit or loss through direct debits or credits to equity reserves.

\item \textbf{Residual Income (Abnormal Earnings) Formulation:}
Residual income (or abnormal earnings, denoted $NI_t^a$) represents the net income generated by the firm in excess of the normal return required on the beginning-of-period book equity capital:
\begin{equation}
NI_t^a = NI_t - r \cdot BV_{t-1}
\end{equation}
By substituting the clean surplus expression for dividends into the dividend-discount model, and assuming the standard regularity transversality condition that the discounted expected book value of equity approaches zero as the horizon approaches infinity ($\lim_{\tau \to \infty} (1+r)^{-\tau} \mathbb{E}_t[BV_{t+\tau}] = 0$), \citet{ohlson1995} proves that the market value of equity can be expressed identically as:
\begin{equation}
P_t = BV_t + \sum_{\tau=1}^{\infty} \frac{\mathbb{E}_t[NI_{t+\tau}^a]}{(1 + r)^\tau} = BV_t + \sum_{\tau=1}^{\infty} \frac{\mathbb{E}_t[NI_{t+\tau} - r \cdot BV_{t+\tau-1}]}{(1 + r)^\tau}
\end{equation}

\item \textbf{Linear Information Dynamics and the Theoretical Valuation Equation:}
To translate this infinite-horizon residual income formulation into an empirical valuation equation containing contemporaneous accounting variables, \citet{ohlson1995} introduces linear information dynamics (LID). Specifically, current abnormal earnings ($NI_t^a$) and other non-accounting information ($v_t$) follow a stationary stochastic autoregressive process:
\begin{align}
NI_{t+1}^a &= \omega NI_t^a + v_t + \epsilon_{1,t+1} \\
v_{t+1} &= \gamma v_t + \epsilon_{2,t+1}
\end{align}
where the persistence parameters satisfy $0 \le \omega < 1$ and $0 \le \gamma < 1$, and $\epsilon_{1,t+1}, \epsilon_{2,t+1}$ are mean-zero disturbance terms.

Under these information dynamics, \citet{ohlson1995} demonstrates that the equilibrium market price of equity ($P_t$) solves as a closed-form linear function:
\begin{equation}
P_t = BV_t + \alpha_1 NI_t^a + \alpha_2 v_t
\end{equation}
where $\alpha_1 = \frac{\omega}{R_f - \omega}$, with $R_f = 1 + r$, and $\alpha_2 = \frac{R_f}{(R_f - \omega)(R_f - \gamma)}$.

Because abnormal earnings are defined with respect to beginning-of-period book value ($NI_t^a = NI_t - r \cdot BV_{t-1}$), the clean surplus relation ($BV_t = BV_{t-1} + NI_t - D_t$, or equivalently $BV_{t-1} = BV_t - NI_t + D_t$) can be substituted directly into the abnormal earnings expression:
\begin{equation}
NI_t^a = NI_t - r(BV_t - NI_t + D_t) = (1 + r)NI_t - r \cdot BV_t - r \cdot D_t
\end{equation}
Substituting this into the valuation equation yields:
\begin{equation}
P_t = BV_t + \alpha_1 \left[(1 + r)NI_t - r \cdot BV_t - r \cdot D_t\right] + \alpha_2 v_t
\end{equation}
Collecting terms on contemporaneous book value ($BV_t$) and net income ($NI_t$) yields the canonical structural form:
\begin{equation}
P_t = (1 - k) BV_t + k (\phi \cdot NI_t - D_t) + \alpha_2 v_t
\end{equation}
where $k = r \cdot \alpha_1 = \frac{r \omega}{R_f - \omega}$ and $\phi = \frac{R_f}{r} = \frac{1+r}{r}$.

Hence, the exact theoretical coefficient on contemporaneous book value of equity is:
\begin{equation}
(1 - k) = 1 - \frac{r \omega}{1 + r - \omega} = \frac{(1 + r)(1 - \omega)}{1 + r - \omega} > 0
\end{equation}
and the theoretical coefficient on contemporaneous net income is:
\begin{equation}
k \cdot \phi = \frac{\omega(1 + r)}{1 + r - \omega} \ge 0 \quad (\text{strictly positive for } \omega > 0)
\end{equation}

When moving from this analytical formulation to empirical testing, the accounting value-relevance literature follows \citet{collins1997} by estimating a linear approximation that omits the unobservable information term ($v_t$) and impounds dividend effects into the constant and error term, scaling variables on a per-share basis:
\begin{equation}
P_{it} = \beta_0 + \beta_1 \text{EPS}_{it} + \beta_2 \text{BVPS}_{it} + \varepsilon_{it}
\end{equation}
where earnings per share ($\text{EPS}_{it}$) serves as the empirical proxy for contemporaneous earnings (with abnormal earning capacity impounded into the structural parameters), and book value per share ($\text{BVPS}_{it}$) serves as the empirical proxy for cumulative net resources.
\end{enumerate}

The Ohlson framework demonstrates that earnings and book value are not arbitrary ad hoc variables selected for regression convenience; rather, they serve as theoretically motivated and complementary inputs in equity pricing. Book value represents the accumulated historical stock of equity net assets, providing an asset-backing floor, while earnings capture current operational flows and abnormal earning power. Second, the theoretical coefficients are direct functions of earnings persistence ($\omega$). As earnings persistence increases, the valuation weight placed on earnings expands. Conversely, if earnings become transitory, noisy, or disrupted by non-recurring events, the persistence parameter $\omega$ declines, leading the market to shift its valuation weight toward balance-sheet net worth.

\subsection{Valuation Dynamics Across Corporate Reorganizations and M\&A}
\label{subsec:ma_valuation}

A central theoretical question in accounting value relevance is whether the valuation multipliers attached to earnings and book values remain constant across major structural corporate reorganizations. In corporate amalgamations and acquisitions, the economic and reporting structure of the enterprise undergoes fundamental reconfiguration.

Investigating the valuation consequences of corporate combinations, \citet{kwon2020} examine changes in the value relevance of accounting information following mergers and acquisitions (M\&A). They demonstrate that structural corporate combinations alter the firm's financial reporting environment through several distinct channels:
\begin{enumerate}[label=\textbf{\arabic*.}]
\item \textbf{Transitory Integration Costs and Earnings Noise:} Post-merger operational integration requires extensive administrative expenditures, harmonization of operating systems, restructuring of branch networks, and rationalization of redundant human resources. These integration expenses generate substantial transitory, non-recurring accounting charges that temporarily suppress reported net income.
\item \textbf{Accounting Harmonization and Portfolio Revaluation:} Merging disparate corporate entities requires harmonizing accounting policies across previously independent balance sheets. In commercial banking, this entails realigning loan classification standards, reassessing non-performing loans (NPLs), unifying impairment provisioning methodologies under NFRS 9 / IFRS 9, and recognizing purchase price allocation adjustments and goodwill.
\item \textbf{Earnings Persistence Impairment:} Because transitional accounting charges and post-merger adjustments are inherently temporary and do not reflect ongoing operational earning capacity, the time-series persistence of reported earnings ($\omega$) declines sharply in the immediate post-merger period.
\item \textbf{Shift Toward Balance-Sheet Anchoring:} Following the residual income logic of \citet{ohlson1995} and the empirical evidence of \citet{collins1997}, when earnings quality and persistence deteriorate due to non-recurring organizational shocks, rational equity investors discount current earnings as a noisy indicator of future performance. Consequently, market participants place greater valuation weight on the stock of net assets---book value per share (BVPS)---which serves as an anchor of fundamental solvency and net worth.
\end{enumerate}

\citet{collins1997} similarly demonstrate across a forty-year empirical panel that the relative value relevance of earnings versus book values shifts systematically across economic regimes: earnings relevance dominates when earnings are persistent and stable, whereas book value relevance increases during periods of structural transition, severe financial stress, and extreme restructuring. These theoretical insights provide the direct intellectual rationale for testing parameter shifts ($\delta_E, \delta_B$) across Nepal's 2023 banking consolidation wave.

\subsection{Neoclassical Valuation Benchmarks: Dividend Capitalization Models}
\label{subsec:dividend_models}

To provide analytical context to accounting valuation, classical financial theory formulated equity pricing around cash dividend distributions:
\begin{itemize}
\item \textbf{The Dividend Irrelevance Proposition:} \citet{miller1961} proved that in a frictionless capital market characterized by zero transaction costs, absence of corporate and personal taxes, symmetric information, and a given corporate investment policy, the dividend payout policy of a firm has no effect on its total market valuation or cost of capital. Firm value is determined solely by its investment decisions and the earning power of its assets. In emerging markets, however, frictionless market assumptions are fundamentally violated by capital market imperfections, asymmetric information, and regulatory mandates.
\item \textbf{Dividend Discount Formulations:} \citet{walter1963} relaxed frictionless assumptions by modeling share price as a function of the internal rate of return on retained earnings ($r_a$) relative to the cost of capital ($k_e$). Walter demonstrated that when $r_a > k_e$ (growth firms), retaining earnings maximizes share value, whereas when $r_a < k_e$, maximizing dividend payout enhances shareholder wealth. \citet{gordon1959} formulated the constant-growth dividend capitalization model:
\begin{equation}
P_0 = \frac{D_1}{k_e - g}
\end{equation}
positing that external investors discount retained earnings more heavily than current dividends due to risk aversion regarding uncertain future capital gains (the ``bird-in-the-hand'' hypothesis).
\end{itemize}

In the Nepalese commercial banking sector, dividend announcements---comprising both cash dividends and capitalized bonus shares---have historically served as vital valuation signals. While the baseline valuation model follows \citet{collins1997} in focusing on earnings and book value, evaluating cash dividends (DPS) in an expanded specification serves as an essential empirical robustness test (Specification R4) to examine whether payout distributions alter the fundamental valuation coefficients.

\subsection{Information Asymmetry, Signaling, and Agency Theories in Banking}
\label{subsec:signaling_agency}

Financial institutions operate in an environment characterized by acute informational friction between bank managers, external equity investors, and regulatory bodies:
\begin{itemize}
\item \textbf{Signaling Theory:} \citet{ross1977}, \citet{bhattacharya1979}, and \citet{john1985} demonstrate that in capital markets plagued by information asymmetry, financial disclosures, dividend announcements, and corporate restructuring decisions serve as costly, credible signals conveying management's private expectations regarding future performance. Because weaker firms cannot afford to mimic costly dividend commitments or stringent audited reporting standards without courting financial distress, published audited accounting statements and declared dividends convey authentic fundamental value.
\item \textbf{Agency Theory:} \citet{jensen1976} and \citet{jensen1986} formulate the free cash flow agency theory, positing that corporate managers have incentives to accumulate resources and pursue empire-building acquisitions rather than returning excess capital to shareholders. In commercial banking, declared cash dividend distributions and transparent audited balance sheets mitigate agency conflicts by constraining managerial discretion and subjecting the institution to regular capital-market monitoring.
\item \textbf{Banking Asset Opacity and Prudential Regulation:} \citet{morgan2002} and \citet{flannery2004} establish that commercial banks are inherently more opaque than industrial corporations. Bank assets consist primarily of unquoted private commercial loans whose true quality and credit risk are shielded from public observation. This severe asset opacity makes statutory auditing, capital adequacy compliance, and standardized regulatory reporting deadlines indispensable for credible market valuation.
\end{itemize}

\subsection{Efficient Market Hypothesis and the Information Content of Disclosures}
\label{subsec:emh_pead}

The Efficient Market Hypothesis (EMH), formulated by \citet{fama1970}, posits that asset prices in an efficient financial market fully reflect all available information. Under the semi-strong form of market efficiency, equity share prices adjust rapidly and unbiasedly to all publicly available information, including audited annual financial accounts, earnings releases, and dividend declarations.

Seminal empirical accounting research by \citet{ball1968} and \citet{beaver1968} established that published accounting earnings possess substantial information content: stock returns systematically adjust in the direction of unexpected earnings announcements. However, subsequent empirical research across global equity markets---most notably documented by \citet{bernard1989}---identified the phenomenon of \textit{post-earnings-announcement drift} (PEAD), wherein equity prices continue to adjust with a temporal lag to published accounting disclosures due to investor underreaction and information processing delays.

In frontier emerging markets such as the Nepal Stock Exchange (NEPSE), market efficiency is further constrained by illiquidity, retail-dominated trading, delayed dissemination of corporate filings, and the absence of a unified electronic corporate announcement timestamp registry. These realities introduce critical timing considerations into value-relevance research. If market prices are observed at fiscal-year-end (Ashadh-end / mid-July), they precede the auditing, board adoption, NRB approval, and public release of annual reports, thereby introducing severe look-ahead bias. Conversely, observing market prices at an objective, standardized window following the statutory reporting deadline ensures the observation date falls after the date by which audited accounts were legally required to be public, allowing sufficient time for market participants to digest and incorporate financial disclosures into secondary market trading prices.

\subsection{Econometric Circularity of Price-Derived Multiples}
\label{subsec:pe_circularity}

A critical theoretical and econometric problem that pervades empirical finance literature in Nepal is the erroneous inclusion of price-derived valuation multiples---specifically the Price-to-Earnings (P/E) ratio---as explanatory independent variables in regressions where Market Price per Share (MPS) is the dependent variable.

The P/E ratio is algebraically defined as:
\begin{equation}
\text{P/E}_{it} = \frac{\text{MPS}_{it}}{\text{EPS}_{it}}
\end{equation}
Specifying an empirical regression model of the form:
\begin{equation}
\text{MPS}_{it} = \beta_0 + \beta_1 \text{EPS}_{it} + \beta_2 \text{P/E}_{it} + \beta_3 \text{BVPS}_{it} + \varepsilon_{it}
\end{equation}
introduces a fatal structural circularity. The dependent variable ($\text{MPS}_{it}$) appears directly in the numerator of the explanatory variable ($\text{P/E}_{it}$). Mathematically:
\begin{equation}
\text{MPS}_{it} = \beta_0 + \beta_1 \text{EPS}_{it} + \beta_2 \left(\frac{\text{MPS}_{it}}{\text{EPS}_{it}}\right) + \beta_3 \text{BVPS}_{it} + \varepsilon_{it}
\end{equation}
Rearranging this equation reveals that $\text{MPS}_{it}$ is being regressed on a function of itself. This specification generates severe mechanical endogeneity, violates the fundamental Gauss-Markov assumption of orthogonality between regressors and error terms ($\mathbb{E}[\text{P/E}_{it} \varepsilon_{it}] \ne 0$), creates artificial multi-collinearity, and inflates the coefficient of determination ($R^2$) through definitional overlap rather than economic explanation. Consequently, theoretical econometric integrity requires that all price-derived multiples be strictly excluded from value-relevance specifications, restricting explanatory regressors to pure accounting fundamentals that are determined independently of secondary market trading prices.

\section{Review of Legal, Regulatory, and Policy Framework in Nepal}
\label{sec:regulatory_framework}

The commercial banking sector and capital market of Nepal operate under a highly structured statutory and prudential regulatory regime:

\subsection{Evolution of Capital Market Regulation: SEC, NEPSE, and SEBON}
\label{subsec:nepse_history}

The institutional development of Nepal's organized capital market commenced with the establishment of the Securities Exchange Centre in 1976 (2033 BS) under the joint sponsorship of Nepal Rastra Bank and the Nepal Industrial Development Corporation (NIDC). The enactment of the Securities Exchange Act, 1983 provided the initial statutory foundation for securities transactions. In 1993, under comprehensive financial sector liberalization programs, the Securities Exchange Centre was restructured into the Nepal Stock Exchange (NEPSE) Limited as the licensed operator of the secondary market, commencing formal floor trading on 13 January 1994 (29 Poush 2050 BS) \citep{nepse2017}. Simultaneously, the Securities Board of Nepal (SEBON) was constituted as the apex capital market regulator to oversee market integrity, protect investor interests, and regulate corporate disclosures \citep{nepse2017}.

\subsection{The Statutory Financial Reporting Framework: BAFIA 2017}
\label{subsec:bafia_framework}

Commercial banks in Nepal operate under the dual regulatory oversight of the Companies Act, 2006 (2063 BS) and the Banks and Financial Institutions Act, 2017 (2073 BS) (BAFIA), alongside prudential directives issued by Nepal Rastra Bank \citep{companiesact2006,bafia2017}.

Regarding financial reporting disclosure, Section 59 of BAFIA, 2017 establishes statutory submission and publication obligations for licensed institutions. Specifically, Section 59(2) delineates two distinct reporting duties:
\begin{enumerate}
\item First, a licensed bank must prepare its balance sheet and profit and loss account and submit them to Nepal Rastra Bank within three months of the completion of each financial year, publishing summary details for the information of the general public within that period as prescribed by NRB; and
\item Second, the bank must publish its audited balance sheet and profit and loss account for the public in the prescribed format within nine months into the subsequent financial year \citep{bafia2017}.
\end{enumerate}

Because commercial bank fiscal years in Nepal terminate at Ashadh-end (mid-July), counting nine months into the subsequent financial year (from Shrawan through Chaitra) terminates at the end of Chaitra (mid-April). Operationalizing the price observation window at the first trading session following the end of Chaitra ($MPS^{RD}$) is the study's institutional operationalization of the statutory nine-month timeline under Section 59(2). In observed banking practice across the sampled institutions, commercial bank boards typically release audited annual highlights and announce dividend recommendations between Ashwin and Poush (October to January), with formal adoption at Annual General Meetings (AGMs) concluded prior to the end-of-Chaitra cutoff. Operationalizing the price observation window at the end of Chaitra provides a reproducible, legally anchored cutoff ensuring that the observation date falls after the statutory deadline by which audited accounts were legally required to be public, eliminating look-ahead bias without relying on fragmented, unverified corporate press-release timestamps.

\subsection{NRB Capital Mandates and the Consolidation Directive}
\label{subsec:nrb_consolidation}

The structural consolidation of Nepalese commercial banks has been driven primarily by explicit central bank policy mandates:
\begin{itemize}
\item \textbf{The 2011 Merger Bylaws:} In 2011 (2068 BS), NRB promulgated the \textit{Bank and Financial Institutions Merger and Acquisition Bylaws}, introducing legal provisions, regulatory exemptions, and procedural guidelines for voluntary mergers \citep{nrb2011}.
\item \textbf{The 2015 Capital Expansion Mandate:} In its Monetary Policy for FY 2015/16 (2072/73 BS), NRB executed an aggressive systemic reform, mandating that all Class 'A' commercial banks quadruple their minimum paid-up equity capital from NPR 2.0 billion to NPR 8.0 billion by the end of FY 2073/74 (mid-July 2017) \citep{nrb2015}. To comply, commercial banks capitalized reserves through extensive bonus share distributions and executed heavy rights share issues, resulting in substantial share capital expansion and subsequent dilution of per-share earnings.
\item \textbf{The 2022--2023 Consolidation Wave:} Recognizing that capital expansion alone had not reduced the excessive number of banking institutions, NRB actively encouraged horizontal ``big mergers'' through moral suasion, regulatory forbearance (such as temporary leniency on Credit-to-Deposit ratios and branch overlapping), and fiscal incentives granted under the Finance Acts.
\end{itemize}

This regulatory drive culminated in the massive consolidation wave of early 2023 (mid-FY 2079/80), wherein major commercial banks executed statutory mergers and acquisitions:
\begin{enumerate}[label=\textbf{\arabic*.}]
\item \textbf{Global IME Bank Limited (GBIME):} Amalgamated with Bank of Kathmandu Limited; unified commercial operations commenced on 9 January 2023 (swap ratio 1:1) \citep{gbime2022agm}.
\item \textbf{Himalayan Bank Limited (HBL):} Acquired Civil Bank Limited; unified commercial operations commenced on 24 February 2023 (swap ratio 100:80.28) \citep{hbl2023report}.
\item \textbf{Kumari Bank Limited (KBL):} Amalgamated with Nepal Credit and Commerce Bank Limited; unified commercial operations commenced on 1 January 2023 (swap ratio 1:1) \citep{kumari2022sgm}.
\item \textbf{Nepal Investment Mega Bank Limited (NIMB):} Formed via the amalgamation of Nepal Investment Bank Limited and Mega Bank Nepal Limited; unified commercial operations commenced on 11 January 2023 (swap ratio 100:90) \citep{nimb2023report}.
\item \textbf{NMB Bank Limited (NMB):} Maintained continuous corporate operations without participating in the early-2023 consolidation wave, serving as the benchmark series. (Note: NMB had previously acquired Kanchan Development Bank Limited in August 2020, adding approximately NPR 0.56 billion to paid-up capital beyond its 13\% bonus share capitalization in FY 2077/78, but remained independent during the 2022--2023 wave).
\end{enumerate}

By mid-October 2023 (Asoj 2080 BS), these amalgamations reduced the total commercial bank population in Nepal to exactly 20 Class 'A' institutions \citep{nrb2023list}.

\section{Empirical Review}
\label{sec:empirical_review}

\subsection{International Empirical Studies}
\label{subsec:intl_empirical}

Empirical accounting research across developed and emerging capital markets provides extensive documentation regarding the association between financial accounting disclosures and equity valuation:

\citet{collins1997} investigated the value relevance of earnings and book values in the United States over a forty-year period (1953--1993) across thousands of firm-year observations. Employing the Ohlson (1995) valuation framework, the authors addressed concerns that the transition from an industrialized to a service- and technology-oriented economy had degraded the relevance of historical accounting numbers. Their empirical findings demonstrated that the combined explanatory power ($R^2$) of earnings and book values had actually increased over the forty-year period. However, they observed a systematic shift in relative explanatory weights: the incremental value relevance of earnings declined, while the incremental value relevance of book value increased, driven by higher incidences of negative earnings, non-recurring charges, and intangible asset intensity.

\citet{kwon2020} investigated whether the value relevance of accounting information shifts following corporate mergers and acquisitions (M\&A) in the United States, using a panel spanning pre- and post-merger observations under SFAS 141(R). Their empirical results demonstrated that M\&A transactions significantly alter the relative valuation roles of earnings and book value. Specifically, during the post-merger integration period, the value relevance of earnings declined due to transitional restructuring expenses, portfolio realignments, and earnings noise. Conversely, the value relevance of book value per share increased significantly, confirming that equity investors rely on balance-sheet net worth as a reliable valuation anchor during corporate reorganization.

In emerging Asian equity markets, \citet{sharma2011} investigated the empirical relationship between accounting fundamentals and equity share prices for firms listed on the Bombay Stock Exchange (BSE 30 index) over 1993/94 to 2007/08. Using pooled OLS and panel regression designs, the study found that earnings per share (EPS), book value per share (BVPS), and dividend per share (DPS) exhibited statistically significant positive associations with share prices, with EPS and BVPS serving as the primary fundamental anchors. \citet{malhotra2013} examined 95 firms listed on the National Stock Exchange of India (NSE 100 index) from 2007 to 2012, demonstrating that book value per share, earnings per share, and price-earnings multiples significantly influenced market prices, whereas firm size and dividend payout exhibited weaker associations. In the Middle East, \citet{almumani2014} evaluated listed commercial banks on the Amman Stock Exchange (ASE) over 2005--2011 using multivariate regression analysis, concluding that EPS, BVPS, and price-earnings ratios explained the predominant share of market price variance, with EPS emerging as the most potent explanatory variable.

\subsection{Foundational Nepalese Empirical Literature (Early to Mid Studies)}
\label{subsec:foundational_nepal}

Empirical investigation into equity valuation determinants in Nepal originated alongside the formalization of NEPSE trading:

\citet{pradhan2003} conducted one of the earliest empirical evaluations of stock market behavior in Nepal, examining listed enterprises over 1992--1997. Employing descriptive and correlational techniques, Pradhan investigated the relative effects of cash dividends and retained earnings on stock prices. The study documented that cash dividends exhibited a stronger, more statistically significant positive association with market share prices than retained earnings, concluding that Nepalese equity investors exhibited strong cash-dividend preference in an environment characterized by institutional opacity and market volatility.

\citet{bhattarai2014} investigated the determinants of share prices across a sample of nine commercial banks (Nabil Bank, Nepal Investment Bank, Standard Chartered Bank Nepal, Everest Bank, Bank of Kathmandu, Machhapuchchhre Bank, Laxmi Bank, Kumari Bank, and Siddhartha Bank) over the period 2006 to 2014. Using pooled OLS regressions, Bhattarai evaluated fiscal-year-end closing market price per share against dividend payout ratio (DPR), dividend yield (DY), earnings per share (EPS), price-earnings (P/E) ratio, and bank size (SIZE). The study reported that EPS and P/E ratios exhibited statistically significant positive associations with market share prices, whereas dividend yield displayed an inverse relationship. However, as noted in Chapter~I, Bhattarai's specification explicitly included P/E as an explanatory regressor alongside MPS, inducing mechanical circularity.

\citet{pradhan2016} examined 14 commercial banks over the period 2002/03 to 2013/14 using a pooled panel regression framework. The authors analyzed the impact of firm size, earnings per share, return on assets, and dividend per share on market equity valuation, alongside macroeconomic controls. Their empirical findings identified bank size and EPS as the most consistent and statistically significant positive determinants of share prices, while macroeconomic variables exhibited secondary explanatory power.

\citet{bajracharya2019} evaluated ten commercial banks listed on NEPSE over the ten-year period from 2006/07 to 2015/16. Using pooled OLS regression models, the authors reported that earnings per share (EPS) and dividend per share (DPS) had a statistically significant positive effect on share prices, whereas return on assets (ROA) and return on equity (ROE) were statistically insignificant in predicting market valuations.

\citet{silwal2019} evaluated ten commercial banks across 100 bank-year observations over FY 2065/66 to 2074/75 (2008/09--2017/18). Employing correlation and multiple regression analysis, the authors examined market price per share in relation to book value per share, return on equity, and P/E ratio. The study concluded that BVPS exerted the largest explanatory influence on stock prices, followed by ROE and P/E ratios, arguing that net asset backing served as the primary benchmark for Nepalese bank investors following the 2015 capital escalation mandate.

\subsection{Empirical Evidence from Recent Nepalese Studies (2021--2024, Pre-Wave Panels)}
\label{subsec:recent_nepal}

Recent empirical research on Nepalese commercial banking valuation documents notable divergence in empirical findings across diverse samples and econometric specifications:

\citet{niroula2021} examined 18 commercial banks over the five-year period FY 2015/16 to 2019/20 (90 bank-year observations). Using pooled regression models, Niroula found that earnings per share had a statistically significant positive impact on market share prices. Conversely, book value per share (BVPS) and return on equity (ROE) exhibited negative coefficients that failed to attain statistical significance, highlighting growing empirical ambiguity regarding the pricing of balance-sheet net worth.

\citet{gyawali2022} evaluated ten commercial banks listed on NEPSE over the five-year span 2017 to 2021 (50 observations). Employing descriptive statistics, Pearson correlation, and multiple regression analysis, Gyawali observed statistically significant positive relationships of share price with EPS and DPS. However, the regression coefficient for return on assets (ROA) was statistically insignificant, suggesting that secondary equity markets discounted accounting return metrics in favor of per-share distributions.

\citet{chhetri2023} investigated 13 commercial banks over an eleven-year panel spanning 2012 to 2022 (143 bank-year observations) using pooled OLS regression. Chhetri regressed market share price on EPS, P/E, BVPS, and ROA. In direct contrast to earlier studies, the empirical results indicated that BVPS, ROA, and P/E were positive and statistically significant determinants of market share prices, whereas earnings per share (EPS) was statistically insignificant. Chhetri argued that institutional investors prioritized balance-sheet asset strength and overall asset productivity rather than reported net earnings per share.

\citet{joshi2023} analyzed seven commercial banks across 70 bank-year observations over FY 2012/13 to 2021/22. Using correlation and linear panel regression techniques, the authors evaluated share price behavior against EPS, DPS, BVPS, and P/E ratios. Their findings indicated that all four explanatory variables exhibited statistically significant positive coefficients, with the P/E ratio and EPS demonstrating the highest explanatory power in predicting NEPSE market prices.

\citet{lamsal2024} evaluated seven commercial banks listed on NEPSE over a multi-year period, examining market price per share in relation to EPS, P/E, DPS, and BVPS. The empirical regression results revealed that EPS and P/E ratios had significant positive effects on share prices, while BVPS yielded an insignificant negative coefficient, reflecting persistent instability in book-value coefficients across sampled institutions.

\citet{giri2024} conducted an empirical investigation into four commercial banks across 40 bank-year observations over the decade FY 2070/71 to 2079/80 (2013/14--2022/23). Employing descriptive and correlational analysis, Giri evaluated share prices against DPS, EPS, P/E, and net worth per share (BVPS). The analysis revealed significant positive correlations between market price and DPS, P/E, and net worth per share, whereas the bivariate correlation with EPS was statistically negligible, reinforcing the argument that cash dividends and net asset backing overshadowed operational earnings among sampled banks.

\citet{gurung2024} evaluated 16 commercial banks across 171 observations over the period 2007 to 2023 using stepwise multiple regression analysis. The study found that earnings per share and book value per share exerted statistically significant positive influences on commercial bank share prices, while dividend per share was statistically insignificant, demonstrating that market valuation was anchored in foundational balance-sheet and earning power rather than payout distributions.

\subsection{Empirical Evidence Spanning the 2023 Consolidation Wave (2024--2026)}
\label{subsec:consolidation_studies}

The most recent body of empirical research has extended observation windows through FY 2023/24 and 2024/25, directly capturing the post-consolidation operating environment:

\citet{pant2024} investigated seven commercial banks across 77 observations over the eleven-year period 2013/14 to 2023/24, spanning the 2023 banking consolidation wave. Using multiple regression models, the authors regressed market price per share on EPS, DPS, BVPS, and P/E. The empirical findings demonstrated that EPS and P/E had statistically significant positive impacts on market share prices. Conversely, the coefficients for dividend per share and book value per share were negative and statistically insignificant, indicating that the rapid asset and equity expansion resulting from banking mergers had weakened conventional book value pricing relationships.

\citet{chaudhari2025} examined ten commercial banks over 2014/15 to 2023/24 (100 bank-year observations). Employing pooled OLS regression, Chaudhari analyzed market price per share against firm-specific variables (EPS, DPS, ROA, bank size) and macroeconomic factors (inflation, interest rates, GDP). The empirical results revealed strong, statistically significant positive coefficients for EPS and DPS. In contrast, return on assets (ROA) displayed a statistically significant \textit{negative} association with market share prices, while bank size and macroeconomic variables failed to attain statistical significance. Chaudhari concluded that equity investors penalized low-margin asset expansion while rewarding per-share earnings generation.

\citet{gurungdahal2025} investigated ten commercial banks listed on NEPSE over the ten-year period FY 2013/14 to 2023/24 using a panel regression framework.\footnote{Gurung and Dahal \citep{gurungdahal2025}, published in \textit{Pravaha} (vol.~31, no.~1, pp.~1--12), is distinct from the comprehensive panel study by Dahal, Gurung, and Ghimire \citep{dahal2025} in the \textit{Interdisciplinary Journal of Management and Social Sciences} (vol.~6, no.~1, pp.~252--261), which examined 209 bank-year observations across all listed commercial banks.} Among the studies reviewed, \citet{gurungdahal2025} provide the closest methodological precedent, because they applied the Hausman specification test and estimated a bank fixed-effects regression model to control for unobserved time-invariant bank heterogeneity. Their empirical results indicated that dividend payout ratio (DPR), P/E ratio, earnings per share, book value per share, and return on assets exerted statistically significant impacts on market share prices. However, Gurung and Dahal estimated static, full-sample coefficients, omitting any formal interaction test of whether valuation parameters shifted before and after the 2023 consolidation wave.

\citet{timalsina2026} evaluated 13 commercial banks over the thirteen-year panel FY 2011/12 to 2023/24 (169 bank-year observations). Using descriptive statistics, Pearson correlation, and panel regression analysis, Timalsina investigated the impact of firm-specific indicators (EPS, DPS, P/E, ROA, size) and macroeconomic factors (GDP, inflation, interest rates) on market share prices. The empirical findings revealed that EPS, DPS, P/E, firm size, and inflation had statistically significant positive associations with share prices, while interest rates exerted a significant negative effect. Return on assets (ROA) and GDP growth showed no statistically significant influence.

In the most contemporary literature, \citet{niraula2026} evaluated 16 commercial banks across 80 observations over FY 2019/20 to 2023/24 using balanced panel OLS regression. The study examined firm-specific determinants, documenting statistically significant positive impacts for Capital Adequacy Ratio (CAR), Cash Reserve Ratio (CRR), and total assets, while ROA was positive but statistically insignificant. \citet{bhandari2026} executed a census study of all 19 NEPSE-listed commercial banks, demonstrating that earnings indicators and dividend payout served as the primary determinants of market prices. In governance and regulatory research, \citet{bhattarai2024} evaluated capital adequacy regulations and commercial bank performance, while \citet{dongol2024} investigated corporate governance mechanisms and financial stability.

Regarding banking consolidation directly, \citet{adhikari2023} conducted an empirical examination of the stock and financial performance of seven commercial banks that underwent mergers between 2013 and 2020. The authors documented that post-merger market share prices experienced substantial downward pressure and reduced dividend payouts, while per-share earnings exhibited high post-merger volatility. In a companion study, \citet{adhikari2023cbr} analyzed the comparative financial performance of merged commercial banks, demonstrating that consolidation failed to generate immediate operational synergies or profitability improvements, as substantial integration friction and restructuring costs offset gains from increased asset scale.

\subsection{Synthesis of Empirical Literature}
\label{subsec:synthesis_matrix}

A critical evaluation of the international and domestic empirical literature reveals four major thematic debates and methodological tensions that characterize commercial bank valuation:

\begin{enumerate}[label=\textbf{\arabic*.}]
\item \textbf{The Earning Power versus Net Asset Value Valuation Anchor:}
A central empirical debate across Nepalese commercial banking literature concerns whether market share prices are primarily driven by current operating profitability (EPS) or balance-sheet equity backing (BVPS). Studies by \citet{bhandari2026}, \citet{chaudhari2025}, \citet{niroula2021}, and \citet{pradhan2016} conclude that EPS serves as the predominant driver of share prices, with book value playing a secondary or statistically insignificant role. Conversely, investigations by \citet{chhetri2023} and \citet{silwal2019} document that BVPS exerts the primary explanatory influence on NEPSE bank equities, with EPS exhibiting negligible or secondary statistical power. This divergence reflects institutional and temporal context: following NRB's 2015 fourfold paid-up capital mandate, equity bases expanded rapidly through bonus share capitalizations, temporarily elevating balance-sheet asset backing as an investor benchmark. However, as capital bases stabilized, operational earning capacity re-emerged as the primary valuation multiple, until the 2022--2023 consolidation wave again restructured balance-sheet equity.

\item \textbf{Econometric Endogeneity and Mechanical Circularity in Valuation Multiples:}
A pervasive methodological flaw across the surveyed literature is the routine inclusion of the Price-to-Earnings (P/E) ratio as an independent explanatory variable in regressions modeling Market Price per Share (MPS) \citep{bhattarai2014,chhetri2023,giri2024,gurungdahal2025,joshi2023,lamsal2024,pant2024,silwal2019,timalsina2026}. Because $\text{P/E} \equiv \text{MPS} / \text{EPS}$, regressing market price on P/E introduces mechanical endogeneity and algebraic simultaneity bias, artificially inflating $R^2$ values and distorting true fundamental regression slopes. Studies that exclude price-derived ratios demonstrate substantially more modest explanatory power and reveal genuine fundamental associations that circular models obscure.

\item \textbf{Temporal Misalignment and Look-Ahead Bias in Price Measurement:}
The vast majority of prior Nepalese valuation studies measure equity share prices at fiscal-year-end (Ashadh-end / mid-July) or employ twelve-month volume-weighted averages \citep{bhattarai2014,pradhan2016,silwal2019}. Under Nepalese corporate reporting reality, audited annual financial accounts are not published at Ashadh-end; they require statutory auditing, NRB prudential clearance, and annual general meeting (AGM) adoption, typically spanning three to nine months into the subsequent fiscal year. Conditioning Ashadh-end market prices on accounting figures that were not publicly disclosed until months later violates the informational sequence of market efficiency, creating severe look-ahead bias and spurious associations.

\item \textbf{Panel Heterogeneity and the Static Pooling Assumption:}
With few exceptions, extant Nepalese banking studies employ pooled ordinary least squares (OLS) regressions across multi-year panels, implicitly assuming parameter homogeneity across disparate institutions and stability across time. While recent work by \citet{gurungdahal2025} introduces bank fixed effects to control for time-invariant bank heterogeneity, existing literature universally estimates static, full-period average coefficients. Consequently, no empirical study has evaluated whether fundamental valuation multiples shifted across the structural regime change induced by Nepal's 2023 commercial banking consolidation wave.
\end{enumerate}

To synthesize the empirical literature directly informing this dissertation's econometric framework, Table~\ref{tab:lit_synthesis} presents a comparative synthesis of primary empirical studies. Selection is restricted to foundational theoretical frameworks and empirical studies whose dependent variable is market price per share (MPS) and which evaluate fundamental accounting indicators across commercial banks.

\begin{table}[!htbp]
\centering
\small
\captionsetup{font={normalsize,bf},justification=justified,singlelinecheck=false}
\caption{Synthesis of Empirical Literature on Accounting Fundamentals and Equity Valuation}
\label{tab:lit_synthesis}
\begin{tabular}{p{0.18\textwidth}p{0.18\textwidth}p{0.16\textwidth}p{0.18\textwidth}p{0.22\textwidth}}
\toprule
\textbf{Study \& Year} & \textbf{Sample \& Period} & \textbf{Methodology} & \textbf{Key Variables} & \textbf{Core Empirical Findings} \\
\midrule
\citet{pradhan2003} & Listed firms, 5 yrs (1992--1997) & Bivariate / Pooled OLS & MPS, DPS, Retained Earnings & Cash dividends show stronger positive price association than retained earnings. \\
\addlinespace
\citet{collins1997} & US Firms, 40 yrs (1953--1993) & Ohlson RIV Panel OLS & Price, EPS, BVPS & Combined relevance increased; book value relevance rose relative to earnings. \\
\addlinespace
\citet{kwon2020} & US M\&A Firms (1998--2014) & Pre/Post M\&A Panel & Price, EPS, BVPS, Post-Merger & Post-merger earnings relevance declines; book value relevance increases as anchor. \\
\addlinespace
\citet{sharma2011} & BSE 30, India, 15 yrs (1993--2008) & Pooled OLS Panel & Price, EPS, BVPS, DPS & EPS, BVPS, and DPS exhibit significant positive associations with equity prices. \\
\addlinespace
\citet{bhattarai2014} & 9 Banks, 72 obs (2006--2014) & Pooled OLS & MPS, DPR, DY, EPS, P/E, Size & EPS and P/E positively associated with MPS; circular P/E inclusion. \\
\addlinespace
\citet{pradhan2016} & 14 Banks, 168 obs (2002--2014) & Pooled Panel OLS & MPS, EPS, ROA, DPS, Size & Bank size and EPS are primary positive determinants; macro controls secondary. \\
\addlinespace
\citet{bajracharya2019} & 10 Banks, 100 obs (2006--2016) & Pooled OLS & MPS, EPS, DPS, ROA, ROE & EPS and DPS have significant positive effects; ROA and ROE statistically insignificant. \\
\addlinespace
\citet{silwal2019} & 10 Banks, 100 obs (2008--2018) & Correlation / Multiple OLS & MPS, BVPS, ROE, P/E & BVPS has largest explanatory power, followed by ROE and P/E; circular P/E. \\
\addlinespace
\citet{niroula2021} & 18 Banks, 90 obs (2015--2020) & Pooled OLS & MPS, EPS, BVPS, ROE & EPS is significantly positive; BVPS and ROE are negative and insignificant. \\
\addlinespace
\citet{gyawali2022} & 10 Banks, 50 obs (2017--2021) & Multiple OLS & MPS, EPS, DPS, ROA & EPS and DPS have significant positive effects; ROA is statistically insignificant. \\
\addlinespace
\citet{chhetri2023} & 13 Banks, 143 obs (2012--2022) & Pooled OLS & MPS, EPS, P/E, BVPS, ROA & BVPS, ROA, P/E positive and significant; EPS insignificant; circular P/E. \\
\addlinespace
\citet{joshi2023} & 7 Banks, 70 obs (2012--2022) & Linear Panel OLS & MPS, EPS, DPS, BVPS, P/E & EPS, DPS, BVPS, and P/E exhibit positive associations; circular P/E. \\
\addlinespace
\citet{giri2024} & 4 Banks, 40 obs (2013--2023) & Correlation Analysis & MPS, DPS, EPS, P/E, Net Worth & MPS correlates with DPS, P/E, net worth; EPS correlation negligible; circular P/E. \\
\addlinespace
\citet{pant2024} & 7 Banks, 77 obs (2013--2024) & Multiple OLS & MPS, EPS, DPS, BVPS, P/E & EPS and P/E positive; DPS and BVPS negative and insignificant; circular P/E. \\
\addlinespace
\citet{chaudhari2025} & 10 Banks, 100 obs (2014--2024) & Pooled OLS & MPS, EPS, DPS, ROA, Size & EPS and DPS strong positive; ROA significantly negative; macro insignificant. \\
\addlinespace
\citet{dahal2025} & All Banks, 209 obs (2013--2024) & Multiple Regression & MPS, P/E, BVPS, DPR, EY, MTB & P/E, BVPS, and dividend rate significantly positive; circular P/E inclusion. \\
\addlinespace
\citet{gurungdahal2025} & 10 Banks, 110 obs (2013--2024) & Bank Fixed Effects & MPS, DPR, P/E, EPS, BVPS, ROA & DPR, EPS, BVPS, ROA significant under bank FE; static pooling; circular P/E. \\
\addlinespace
\citet{timalsina2026} & 13 Banks, 169 obs (2011--2024) & Panel Regression & MPS, EPS, DPS, P/E, ROA, Size & EPS, DPS, P/E, size, inflation positive and significant; circular P/E. \\
\addlinespace
\citet{bhandari2026} & 19 Banks, Census (2015--2024) & Multiple OLS & MPS, EPS, DPS, DPR, P/E & EPS and dividend indicators primary determinants; circular P/E. \\
\addlinespace
\citet{adhikari2023} & 7 Merged Banks (2013--2020) & Pre/Post Comparative & MPS, EPS, DPS, Net Profit & Substantial post-merger price drops and dividend cuts; slow operational synergy. \\
\bottomrule
\end{tabular}
\end{table}

\section{Research Gap}
\label{sec:research_gap}

A review of the theoretical foundations, regulatory frameworks, and empirical literature reveals five key conceptual, methodological, and econometric gaps that motivate this dissertation:

\subsection{The Consolidation Regime-Shift Gap}
Almost without exception, existing Nepalese banking studies employ pooled cross-sectional or static panel regression designs that estimate sample-average, full-period relationships \citep{bhattarai2014,chaudhari2025,chhetri2023,gurungdahal2025,timalsina2026}. These models implicitly assume parameter constancy across time. However, early 2023 witnessed the most concentrated commercial banking consolidation wave in Nepalese history, altering reporting boundaries, asset bases, and capital structures across four major merged institutions. Following the theoretical framework of \citet{kwon2020} and \citet{collins1997}, major corporate amalgamations impair earnings persistence through integration costs and loan harmonization, suggesting a possible shift in valuation weight toward balance-sheet equity. Despite this theoretical precedent, to the best of the researcher's knowledge and within the extensive corpus of studies reviewed, none has formulated an explicit regime-interaction econometric model to test (in a two-tailed framework) whether the valuation coefficients of EPS ($\beta_E$) and BVPS ($\beta_B$) differed between the pre-consolidation (FY 2076/77--2078/79) and post-consolidation (FY 2079/80--2081/82) periods.

\subsection{The Specification and Endogeneity Gap (Circularity of P/E Multiples)}
A pervasive practice in the Nepalese empirical literature is the routine inclusion of the Price-to-Earnings (P/E) ratio as an independent regressor in models predicting Market Price per Share (MPS) \citep{bhattarai2014,chhetri2023,giri2024,gurungdahal2025,joshi2023,lamsal2024,pant2024,silwal2019,timalsina2026}. As mathematically demonstrated in Section~\ref{subsec:pe_circularity}, because $\text{P/E} = \text{MPS} / \text{EPS}$, regressing MPS on P/E introduces a circular algebraic dependency, embedding the dependent variable directly inside the regressor. This induces severe mechanical endogeneity, inflates $R^2$ values spuriously, and invalidates statistical inference. A defensible research design must eliminate price-derived multiples and restrict explanatory variables to fundamental accounting measures determined independently of secondary market share prices.

\subsection{The Temporal Alignment and Look-Ahead Bias Gap}
Prior studies in Nepal frequently measure share prices at fiscal-year-end (Ashadh-end / mid-July) or use annual trading averages \citep{bhattarai2014,pradhan2016,silwal2019}. Because audited annual financial statements are not prepared, audited, and published until months after Ashadh-end, pairing Ashadh-end prices with audited numbers creates severe look-ahead bias, violating temporal causality. Furthermore, attempting to reconstruct individual corporate announcement dates across 30 bank-years is prone to measurement error, informal leakage, and archival opacity. Existing literature lacks a standardized, institutionally grounded price-measurement rule. This study operationalizes Section 59(2) of BAFIA, 2017 at the end of Chaitra (mid-April), observing market prices at the first subsequent trading session ($MPS^{RD}$), thereby ensuring that the observation date falls after the statutory deadline by which audited accounts were legally required to be public while maintaining a reproducible observation threshold.

\subsection{The Entity-Timing Transition Gap}
None of the empirical studies reviewed addresses the institutional transition timing of merged banks. For fiscal year 2078/79, annual financial reports represent pre-merger standalone operations, whereas by mid-April 2023, secondary market trading reflected the unified post-merger entities (mergers completed in January--February 2023). This creates an entity-timing mismatch between pre-merger accounting fundamentals and post-consolidation market valuations. Prior studies ignore this structural shock, risking biased parameter estimation. This study explicitly addresses entity transition timing through dedicated sensitivity exclusions.

\subsection{The Econometric and Small-Cluster Inference Gap}
Prior empirical research on Nepalese commercial banking routinely overlooks finite-sample degrees of freedom and the econometric hazards of small-cluster panel inference. In fixed-effects panel models with interactions, parameter estimation rapidly consumes degrees of freedom. Moreover, standard cluster-robust standard errors suffer from severe downward bias when the number of clusters is small ($G = 5$) \citep{cameron2008clusters}. Published Nepalese studies frequently report unadjusted OLS or naive clustered standard errors without acknowledging finite-sample fragility. This dissertation addresses this gap by implementing MacKinnon-White HC3 heteroskedasticity-robust standard errors, disclosing exact residual degrees of freedom ($28 - 14 = 14$ df), and framing findings as observed statistical associations rather than causal proofs.

By synthesizing the \citet{ohlson1995} residual income valuation model with the post-merger value-relevance framework of \citet{kwon2020}, operationalizing the standardized BAFIA reporting deadline ($MPS^{RD}$), excluding circular P/E multiples, and executing a rigorous two-way fixed-effects panel model with eight robustness checks, this dissertation is designed to address each of the identified gaps in the literature.

